\section{Evaluation}
\label{sec:eval}

The evaluation answers six research questions.
\textbf{RQ1 (Varying the compilation):} what does varying the compilation find that varying the program does not?
\textbf{RQ2 (Program dimension):} what do metadata-derived program-side factors find?
\textbf{RQ3 (Oracle):} which findings require an observable other than the output value?
\textbf{RQ4 (Attribution):} does layered execution name the stage that the developers eventually fix?
\textbf{RQ5 (Comparison):} how does the one-factor matrix compare with existing generators and with a cold-only
comparison on a controlled benchmark?
\textbf{RQ6 (Real-world defects):} what did the approach find in current PyTorch builds, and how did the
maintainers respond?

\subsection{Experimental Setup}
\label{sec:setup}

\paragraph{Implementation.} \tool is implemented in about 25{,}000 lines of Python on top of the standard
\code{ast} module and PyTorch's public and testing interfaces (\code{backend=}, Dynamo counters, OpInfo,
\code{TorchDispatchMode}, fake tensors, \code{torch.export}, AOTInductor). Every work item runs in a child
process with resumable per-item records. Disk caches are disabled unless they are the object of the test.

\paragraph{Subjects and environments.} The subject is PyTorch~2.14.0, released during the study, and nightly
builds of the development branch (2.15.0.dev, September 2026). Experiments ran on Windows~11 with MSVC and
Python~3.14 (CPU), on Linux with GCC and Python~3.11/3.12 (CPU), and on Linux with a Tesla~T4 GPU (CUDA~13,
Triton). PyTorch~2.10.0 on the same Linux image serves as the older version for the historical benchmark.

\paragraph{Benchmarks and seeds.} (i)~A \emph{controlled fault benchmark} of 19 faults injected through
compiler backends that wrap the real pipeline and corrupt one decision (a generated value, a shape boundary or
vectorization tail, a dtype, a branch, a dropped mutation or a view turned into a copy, an under-specialized
guard, a stale or memoized cache entry), each attached to one of 29 corpus programs; the unmodified compiler on
the same programs is the fixed version. (ii)~A \emph{historical benchmark} of 193 reproducers mined from the
PyTorch issue tracker, run on PyTorch~2.10 and~2.14. (iii)~\emph{Seeds}: the 29-program corpus, the
274-program Python corpus, OpInfo samples for 703 operator entries, NNSmith-generated models, the mined
reproducers, and 73 programs built from 14 torchvision architectures.

\paragraph{Baselines.} Random context generation (B1) draws values for every factor family without consulting
the program; NNSmith~\cite{DBLP:conf/asplos/LiuLRTLPZ23} (B3) is the generator closest to our setting that supports
PyTorch~2. Both run inside our harness with the same budget, compiler, and machine, with and without our
oracles. TorchProbe~\cite{DBLP:conf/aplas/SuGPS23}, FreeFuzz~\cite{DBLP:conf/icse/WeiDYZ22}, and TitanFuzz~\cite{DBLP:conf/issta/DengXPY023} could not be
executed against PyTorch~2.14 in our environment (Section~\ref{sec:threats}).

\subsection{RQ1: Varying the Compilation}
\label{sec:rq1}

\paragraph{One factor on a covered corpus.} The Python corpus allows the program set, the inputs, and the
oracle to be held fixed while only the path changes. Program-only
differential testing had run the 185 programs of the corpus at that time under three configurations,
\code{backend="eager"}, \code{dynamic=True}, and \code{aot\_eager}: 19 divergences, all of them known. Adding one path factor, a graph break after every
statement, raised the count from 19 to 25 in the same process with the same oracle: eight new divergences (four
programs, two calls each) and two that disappeared. The new ones belong to two defects, both silent: an \code{OrderedDict.move\_to\_end} on a dictionary that enters
the compiled function is dropped by the side-effect replay (C41 in Table~\ref{tab:bugs}), and a generator
that is alive across the break is reconstructed as an exhausted tuple iterator, so that \code{.send()} and
\code{.throw()} raise \code{AttributeError} and \code{StopIteration.value} is \code{None} (C42). Both
reproduce on Linux with Python~3.12 and on the nightly build. Neither needs the break once the mechanism is
known: the first needs only a dictionary passed as an argument, the second only a
generator that is returned. The break exposed them because it turns a local object created inside the frame into an input of the
resume function. Replacing the break by a \code{disable}-wrapped call or by a
\code{print} produced the same eight divergences; a pure Python side effect without a break produced none;
breaks inside loop, conditional, and \code{try} bodies (816 insertion points) added nothing to the top-level
variant (660). Under the Inductor backend the same factor additionally exposed the identity family of C8
(\code{x * 1} after a break returns the input object), which had been found earlier through a program-side
factor.

\paragraph{Programs added after the findings.} After C41 and C42, 57 programs were added. Most exercise the
mechanisms involved (context managers, the generator protocol driven by the consumer, ordered containers passed as
inputs, lazy built-ins with side effects); a smaller group covers exception flow around tensor work. Under the
break factor this batch added eight divergences, all members of C41 and C42, and two on a
\code{functools.lru\_cache} wrapper whose inlining is documented behavior. On the unmodified path it exposed
C44, an \code{IndexError} raised by a tensor operator that bypasses the program's own \code{except} clause, a
mechanism unrelated to the break. The 33 programs added after C44 were written for that mechanism: they showed
that every handler shape (\code{except Exception}, \code{contextlib.suppress}, a swallowing
\code{\_\_exit\_\_}, nested calls, loops, generators) is bypassed and that \code{NotImplementedError} takes the
same route (C45), while \code{TypeError}, \code{ZeroDivisionError}, assertions, and \code{except*} groups reach
their handlers. Twenty-two of the 33 programs diverged, which delimits the scope of C44 without adding a defect.

\paragraph{Other path factors.} The forwarding dispatch mode compared 955 differentiable operator--dtype pairs
and their gradients; three operators differed (\code{linalg.pinv}, \code{matrix\_sqrth}, \code{linalg.lstsq},
all complex), and confirmation without \tcompile attributed them to conjugate views that are not materialized
under any dispatch mode (C32, C39), which also corrected the attribution of B1 from AOTAutograd to the dispatch
machinery (C33). The artifact route compared 520 operators through export, serialization, and AOTInductor: the
values agreed with \tcompile for every operator; the one route-specific finding is a process crash when the
loaded runner is destroyed immediately after a run on Windows (C46). Compile composition, autocast, higher-order operators, custom operators, and CUDA-graph output reuse found
nothing beyond documented behavior (Table~\ref{tab:negative}).

\rqanswer{RQ1}{Of the 58 findings attributable to the compile pipeline, 20 came from a path-side factor. On a
185-program corpus that program-only testing under three configurations had already covered, one path factor
exposed two new defects and one attribution correction; the 90 programs added afterwards found a third on the
unmodified path and delimited it; the dispatch-mode
and artifact-route factors each contributed a family that program-side variation does not reach.}

\begin{table}[t]
\caption{Factor families that found no new inconsistency on PyTorch~2.14 (P = program side, Q = path side).}
\label{tab:negative}
\scriptsize
\setlength{\tabcolsep}{4pt}
\begin{tabular}{@{}lp{5.4cm}p{3.5cm}p{3.4cm}@{}}
\toprule
Side & Factor family & Scale & Outcome \\
\midrule
Q & in-memory guards (dtype, layout, \code{requires\_grad}, flags) & 26 single-factor switches & always recompiled; cold = warm \\
Q & on-disk caches across processes & 29 baked-in values and global switches & 27 misses, 2 correct hits, 0 stale \\
Q & Inductor code-generation switches & 593 ops $\times$ 20 configurations; 68 fused programs $\times$ 15 & 0 violations \\
Q & compile composition (\code{compile(compile(f))}, nested calls, \code{disable}, warm-up in \code{no\_grad}, two backends) & 14 programs $\times$ 25 compositions & 21 differences, all documented behavior \\
Q & autocast, bfloat16 and float16 & 530 ops $\times$ 2 dtypes & 0 dtype or value differences \\
Q & higher-order operators (\code{cond}, \code{while\_loop}, \code{scan}, checkpoint) & 47 programs $\times$ 3 backends & every difference is a raised error \\
Q & custom operators and \code{autograd.Function} & 23 programs $\times$ 3 backends & 0 silent (one dispatch-mode member of C32) \\
Q & export serialization and AOTInductor values & 520 ops & values identical to \tcompile \\
Q & CUDA-graph output reuse (\code{reduce-overhead}) & 19 access patterns (T4) & 13 protected by the documented error, 6 unchanged, 0 stale \\
Q & graph breaks inside loop / conditional / \code{try} bodies; side-effect replay without a break & 816 insertion points; 471 markers & nothing beyond the top-level variant; effects replayed once \\
P & \code{vmap} inside the compiled region; forward-mode AD with complex inputs & 678 ops; 506 ops & 0 \\
P & binding form of scalar arguments & 386 ops, 10{,}651 compilations & 0 new \\
P & small real models (torchvision), CPU and GPU & 131 programs & 0 \\
P & binary functions and reductions against a 30-digit reference & 23 functions, 18 reductions & 9 flagged, all previously reported \\
\bottomrule
\end{tabular}
\end{table}


\subsection{RQ2: The Program Dimension}
\label{sec:rq2}

Metadata-derived factors keep the operator call fixed, so a sweep costs one compilation per pair. The unsupported-dtype analyzer ran
1{,}345 operator--dtype pairs on CPU: in 418 pairs of 134 operators the compiled program returned a value where
the eager kernel raises \code{NotImplementedError} (C37), in five mechanisms (integer operators promoted to
float, Boolean inputs, bitwise operators on floating-point bit patterns, half-precision special functions,
complex normalizations); on the T4 the same factor gave 383 pairs, 73 of them CUDA-specific (C38). The
boundary-value analyzer found nothing at the extreme values of the first version of its table and four
functions after values next to singularities were added (\code{nextafter(1, 0)}, $1 \pm \epsilon$,
$\sqrt{\mathit{max}}$): \code{erfinv} loses four digits near $\pm 1$ (C35), the vectorized float64
\code{acosh} overflows above $1.3\times10^{154}$ (C36), and the vectorized float32 \code{erf} returns 0 for
$|x| \le 10^{-7}$ (C40). Comparing 44 unary functions against a 30-digit reference flagged six function--dtype
pairs, of which \code{erf} and \code{atanh} were new; the same comparison for 23 binary functions and 18
reductions flagged nine, all of them families already reported (three by us); the analyzer thus re-finds known
defects. Two lessons shaped the analyzer: probe tensors must be at least the vector width long, since a
shorter probe tests only the scalar tail of the generated loop, and eager execution is not a sufficient
accuracy reference, since eager and compiled \code{erf} agree to $10^{-7}$ absolutely while the compiled
result has no correct digit below $10^{-7}$.

\rqanswer{RQ2}{Program-side factors derived from operator metadata account for 38 of the 58 pipeline findings;
one factor, casting to a dtype the eager kernel rejects, exposed 134 operators, and a multi-precision reference
exposed accuracy losses that an eager-versus-compiled comparison cannot see.}

\subsection{RQ3: Contribution of the Oracle}
\label{sec:rq3}

Table~\ref{tab:bugs} names the observable that detected each issue. Of the 58
findings attributable to the pipeline, 19 are wrong values that a comparison of outputs against eager execution
detects. The remaining 39 need another observable: 13 are failures of programs that eager execution accepts
(internal assertions, C++ compilation errors, lowering exceptions), 9 are lost validations in which the compiled
program returns where eager execution raises, 4 are dtype changes, 3 are identity changes (an output that
\emph{is} the input), 2 are wrong gradients with a correct forward value, 2 are dropped side effects, 2 are
accuracy losses visible only against a multi-precision reference, 2 are process crashes, and 1 each are a
memory-layout and a device change. The consequences are not local to the reported value: the identity change of C8 turns a later in-place
update into corruption of the input, the dropped \code{move\_to\_end} of C41 makes an LRU cache evict the
wrong entry, and the bypassed handler of C44 disables a fallback that works in eager mode.

\rqanswer{RQ3}{Thirty-nine of the 58 pipeline findings are invisible to an output-value comparison; exception
parity in both directions accounts for 22 of them, and identity, state, gradient, precision, process, layout,
and device observables for the rest.}

\subsection{RQ4: Attribution}
\label{sec:rq4}

On the controlled benchmark, the layer reported by localization matches the layer of the injected fault in all
19 cases. For real defects the ground truth is the fix. Thirteen PyTorch findings have a fix landed or under
review; in all 13 the files changed by the fix lie in the layer that localization named (TorchDynamo for B2,
B17, B26, C11; Inductor lowerings and decompositions for B8, B21, B23, B24; the reference layer for B14;
AOTAutograd for C1; export packaging for B20; the shared eager kernel for B16). Ten findings were attributed outside the pipeline: three eager kernels reached through the compiled
path (B16, C2, C4), three in \code{torch.fx} code generation and serialization (C9, C24, C27), and four in the
dispatch machinery (B1 after correction, C32, C33, C39). They are reported and counted separately, because a
differential oracle detects a divergence between two executions and cannot by itself decide which side is
wrong.

\rqanswer{RQ4}{Localization named the layer of all 19 injected faults and the layer of all 13 fixed real
defects; the same procedure attributed 10 findings to components outside the compile pipeline, which the paper
counts apart.}

\subsection{RQ5: Comparison on the Controlled Benchmark}
\label{sec:rq5}

%% Figure: scatter of faults found vs tests executed; marker area encodes test efficiency.
\begin{figure}[t]
\centering
\begin{tikzpicture}
\begin{axis}[
    tccaxis, width=0.78\linewidth, height=4.3cm,
    xmin=100, xmax=560, ymin=0, ymax=21,
    xlabel={tests executed (budget: 12 per program)}, ylabel={injected faults found (of 19)},
    ytick={0,5,10,15,19},
    xmajorgrids,
    clip=false,
]
  %% x = tests, y = faults, radius ~ sqrt(test efficiency)
  \addplot[only marks, mark=*, mark size=6.3pt, draw=tolred!80!black, fill=tolred!35] coordinates {(235,19)};
  \addplot[only marks, mark=*, mark size=4.5pt, draw=tolblue!80!black, fill=tolblue!30] coordinates {(355,14)};
  \addplot[only marks, mark=*, mark size=3.2pt, draw=tolblue!80!black, fill=tolblue!12] coordinates {(464.7,11.33)};
  \addplot[only marks, mark=*, mark size=4.0pt, draw=tolteal!80!black, fill=tolteal!12] coordinates {(152,6)};
  \addplot[only marks, mark=*, mark size=2.7pt, draw=tolteal!80!black, fill=tolteal!35] coordinates {(496,9)};
  \node[anchor=west, font=\scriptsize] at (axis cs:250,19.3) {\tool (0.098)};
  \node[anchor=west, font=\scriptsize] at (axis cs:367,14.4) {B1 random, all oracles (0.051)};
  \node[anchor=east, font=\scriptsize] at (axis cs:455,11.0) {B1 random, output oracle (0.025)};
  \node[anchor=west, font=\scriptsize] at (axis cs:163,5.6) {B3 NNSmith, native inputs (0.040)};
  \node[anchor=east, font=\scriptsize] at (axis cs:488,8.4) {B3 NNSmith + \tool contexts (0.018)};
  \draw[black!40, densely dashed] (axis cs:100,19) -- (axis cs:560,19);
\end{axis}
\end{tikzpicture}
\caption{RQ2: faults found against tests executed under the same harness, oracles, and budget. Marker area is
proportional to test efficiency (unique failure clusters per executed test, in parentheses). B1 values are means
over three seeds.}
\label{fig:rq2}
\end{figure}


Figure~\ref{fig:rq2} plots, for a budget of 12 tests per program, the injected faults found against the tests
executed. Random generation with the output-only oracle that is customary in differential testing finds 11.3 of
19 faults (mean of three seeds) in 465 tests, 19\% of which are invalid; giving it all seven observables raises
it to 14, so one third of the gap is due to the oracle and two thirds to the choice of factors. NNSmith's
models expose no Python-level factors, so even with our contexts added detection stays at 9. \tool finds all 19 faults in 235 tests, with 16 invalid
contexts, and raises no alarm on the unmodified compiler. An ablation that enables the analysis layers
cumulatively (Figure~\ref{fig:rq3}) raises detection from 7 to 19 while the number of invalid tests falls from
109 to 16; alias facts and context sequences account for 9 of the 12 additional detections. Six faults require
a warm artifact: a cold-only comparison finds none of them in 148 tests, random sequences find 5, and
\scs-guided sequences find all 6 with 40\% fewer tests (Figure~\ref{fig:rq4}). Against the real compiler the
same sequences and 29 cross-process cache cases revealed no stale artifact. On the 193 mined reproducers
(Figure~\ref{fig:rq1-history}), detection of open issues stays at 41\% across PyTorch~2.10 and~2.14 while
detection of closed issues drops from 41\% to 13\%, which yields 35 verified buggy/fixed pairs.

%% Figure: lollipop chart for cache/specialization faults + test counts as annotated bars.
\begin{figure}[t]
\centering
\begin{tikzpicture}
\begin{groupplot}[
    group style={group size=2 by 1, horizontal sep=0.5cm, y descriptions at=edge left},
    tccaxis, height=3.4cm,
    ytick={1,2,3},
    yticklabels={ordinary eager vs.\ compiled, random sequences, \scs-guided sequences},
    ymin=0.4, ymax=3.6, ymajorgrids=false, xmajorgrids,
]
  \nextgroupplot[width=0.40\linewidth, xmin=0, xmax=6.6, xtick={0,2,4,6}, xlabel={cache faults detected (of 6)}]
    \addplot[xcomb, very thick, tolblue, mark=*, mark size=2.6pt] coordinates {(0,1) (5,2) (6,3)};
    \node[font=\scriptsize, anchor=west] at (axis cs:0.15,1) {0};
    \node[font=\scriptsize, anchor=east] at (axis cs:4.8,2.32) {5};
    \node[font=\scriptsize, anchor=east] at (axis cs:5.8,3.32) {6};
  \nextgroupplot[width=0.40\linewidth, xmin=0, xmax=190, xlabel={tests executed (mean tests to trigger)}]
    \addplot[xbar, bar width=6pt, fill=tolorange!60, draw=tolorange!80!black] coordinates {(148,1) (155,2) (93,3)};
    \node[font=\scriptsize, anchor=west] at (axis cs:149,1) {(--)};
    \node[font=\scriptsize, anchor=west] at (axis cs:156,2) {(6.5)};
    \node[font=\scriptsize, anchor=west] at (axis cs:94,3) {(4.8)};
\end{groupplot}
\end{tikzpicture}
\caption{RQ4: the six faults that need a warm artifact. A cold-only comparison cannot observe them; \scs guidance
finds all six with fewer tests than random sequences.}
\label{fig:rq4}
\end{figure}

%% Figure: ablation -- faults detected (step line) and valid/invalid tests (stacked bars).
\begin{figure}[t]
\centering
\begin{tikzpicture}
\begin{groupplot}[
    group style={group size=2 by 1, horizontal sep=1.5cm},
    tccaxis, width=0.49\linewidth, height=3.9cm,
    symbolic x coords={A0,A1,A2,A3,A4,A5,A6,FULL},
    xtick=data, x tick label style={font=\scriptsize},
]
  \nextgroupplot[ymin=0, ymax=23, ylabel={faults detected (of 19)}, ytick={0,5,10,15,19}]
    \addplot[const plot mark mid, very thick, tolred, mark=*, mark size=1.8pt]
      coordinates {(A0,7) (A1,8) (A2,8) (A3,10) (A4,13) (A5,13) (A6,19) (FULL,19)};
    \node[font=\tiny, anchor=south east] at (axis cs:A3,10.3) {+shape/index};
    \node[font=\tiny, anchor=south east] at (axis cs:A4,13.3) {+alias facts};
    \node[font=\tiny, anchor=south east] at (axis cs:A6,19.3) {+sequences};
  \nextgroupplot[ybar stacked, bar width=7pt, ymin=0, ymax=460, ylabel={tests executed},
                 legend pos=north east, legend columns=2]
    \addplot[fill=tolblue!45, draw=tolblue!80!black] coordinates
      {(A0,229) (A1,279) (A2,287) (A3,274) (A4,193) (A5,228) (A6,219) (FULL,219)};
    \addplot[fill=tolorange!70, draw=tolorange!80!black] coordinates
      {(A0,109) (A1,71) (A2,71) (A3,69) (A4,69) (A5,16) (A6,16) (FULL,16)};
    \legend{valid, invalid}
\end{groupplot}
\end{tikzpicture}
\caption{RQ3: cumulative ablation of the analysis layers (A0: dynamic random contexts; A1: tensor metadata;
A2: flow-sensitive dependencies; A3: shape--index relations; A4: alias and mutation facts; A5: \scs pruning;
A6: context sequences). Detection rises while the number of tests, and in particular of invalid tests, falls.}
\label{fig:rq3}
\end{figure}

%% Figure: slope chart of detection rates on mined issue reproducers across versions.
\begin{figure}[t]
\centering
\begin{tikzpicture}
\begin{axis}[
    tccaxis, width=0.8\linewidth, height=4.0cm,
    xmin=-0.75, xmax=3.7, ymin=0, ymax=60,
    xtick={1,2}, xticklabels={PyTorch 2.10, PyTorch 2.14},
    ylabel={issues detected (\%)},
    clip=false,
]
  \addplot[tolorange, very thick, mark=*, mark size=2.2pt] coordinates {(1,45.7) (2,41.4)};
  \addplot[tolblue, very thick, mark=square*, mark size=2.2pt] coordinates {(1,40.7) (2,13.0)};
  \node[anchor=east, font=\scriptsize, tolorange!80!black] at (axis cs:0.9,48.5) {open issues: 32/70};
  \node[anchor=east, font=\scriptsize, tolblue!80!black] at (axis cs:0.9,37.5) {closed issues: 50/123};
  \node[anchor=west, font=\scriptsize, tolorange!80!black] at (axis cs:2.07,41.4) {29/70};
  \node[anchor=west, font=\scriptsize, tolblue!80!black] at (axis cs:2.07,13.0) {16/123};
  \draw[black!55, decorate, decoration={brace, amplitude=3pt, mirror}] (axis cs:2.55,14) -- (axis cs:2.55,40);
  \node[anchor=west, font=\scriptsize, black!75, align=left] at (axis cs:2.66,27) {35 verified\\buggy/fixed pairs};
\end{axis}
\end{tikzpicture}
\caption{RQ1 on 193 reproducers mined from the PyTorch issue tracker, executed on the same Linux CPU image with
two PyTorch versions. Open issues remain detectable; closed issues stop being detected once their fix ships.}
\label{fig:rq1-history}
\end{figure}


\rqanswer{RQ5}{Under identical oracles and budget, the one-factor matrix finds 19 faults against 14 for random
generation and 9 for NNSmith with half the tests; on the real compiler, guards and caches were consistent, and
the defects of RQ6 were reached through the paths behind them.}

\subsection{RQ6: Defects in Current PyTorch Builds}
\label{sec:rq6}

\paragraph{Reports.} The sweeps on PyTorch~2.14 and nightly builds have produced 73 items to date
(Table~\ref{tab:bugs} lists the issues): 55 new issues, 15 comments that add members, platforms, or a
corrected root cause to existing issues, and 3 reports to the vendor of the C++ compiler that TorchInductor
uses on Windows, whose optimizer folds vector intrinsics that PyTorch's own vectorized kernels emit. Every
candidate passed the confirmation steps of Section~\ref{sec:phase3}, was reproduced on PyTorch~2.14 and on a
nightly build, on Windows and Linux (and on CUDA where applicable), and was searched in the issue tracker by
root-cause keywords before filing. By maintainer response, 14 are fixed or have a fix under review, 10 are
confirmed by a maintainer or reproduced by another contributor, 48 await a human response, and 1 was judged
expected behavior (a mutated captured constant that export accepts and AOTInductor rejects). By root cause, 58
lie in the compile pipeline, 10 outside it, and 5 in other compilers (JAX, Numba, and the three MSVC reports). Of the 58,
42 are silent, 14 are failures of programs that eager execution accepts, and 2 terminate the process; they
span TorchDynamo (18), TorchInductor (29), AOTAutograd (4), decompositions and meta kernels (5), and export and
AOTInductor (3). Every change to \tcompile passes a continuous-integration suite that already runs OpInfo-based compiler
tests, and the tracker is triaged daily by dedicated teams.

\begin{table}[t]
\caption{The 55 issues filed against PyTorch and two other compilers (identifiers anonymized), in order of filing.
\emph{Layer}: Dy = TorchDynamo, AOT = AOTAutograd / autograd tracing, De = decompositions and meta kernels,
In = TorchInductor, Ex = export / AOTInductor, Eg = eager kernel shared by the compiled path, Fx = \code{torch.fx},
DM = dispatch machinery, Ext = another compiler. \emph{Sym.}: S = silent inconsistency, F = failure of a program
that eager execution accepts, C = process crash. \emph{Fac.}: P = program-side factor, Q = path-side factor.
\emph{Obs.}: observable that detected it (V value, D metadata, I identity, E exception, S state, G gradient,
R precision reference, C process). \emph{St.}: F fixed or fix under review, C confirmed, P pending, R rejected.}
\label{tab:bugs}
\scriptsize
\setlength{\tabcolsep}{3pt}
\begin{tabular}{@{}rp{8.2cm}llllll@{}}
\toprule
ID & Inconsistency (eager vs.\ compiled) & Layer & Sym. & Fac. & Obs. & St. \\
\midrule
B1  & \code{linalg.pinv} / \code{matrix\_sqrth}: wrong gradient for complex inputs & AOT & S & P & G & P \\
B2  & \code{random.shuffle} / \code{sample} return the same result on every call & Dy & S & Q & V & F \\
B3  & \code{jvp} of \code{ldexp} terminates the interpreter under \ET & AOT & C & Q & C & P \\
B4  & arg-reduction index emitted with Python \code{**} in C++; compilation fails & In & F & Q & E & C \\
B5  & \code{new\_zeros(device="cpu")} returns a CUDA tensor & In & S & Q & D & C \\
B6  & \code{std} / \code{var}: NaN for $10^{30}$, zero value \emph{and zero gradient} for $10^{-30}$ & In & S & P & V & P \\
B7  & \code{jvp} of \code{quantile} / \code{addr}: internal assertion on every backend & Dy & F & Q & E & P \\
B8  & \code{ldexp} result dtype wrong in both directions & In & S & P & D & F \\
B9  & compiled autograd: \code{interpolate} backward hits an internal assertion & AOT & F & Q & E & P \\
B10 & f-string format specifier on a symbolic size: internal Dynamo error & Dy & F & Q & E & C \\
B11 & compiled \code{SGD.step}, \code{foreach=True}, complex parameters: lowering assertion & In & F & P & E & C \\
B12 & \code{interpolate} on an empty spatial dim returns NaNs (out-of-bounds read); eager raises & De & S & P & E & C \\
B13 & mutated captured tensor: export accepts, decomposition rejects, AOTInductor asserts & Ex & F & Q & E & R \\
B14 & \code{vector\_norm(ord=inf)} on an empty batch fails only for negative \code{dim} & De & F & P & E & F \\
B15 & \code{binary\_cross\_entropy} / \code{huber\_loss}: float32 result for half-precision inputs & In & S & P & D & P \\
B16 & \code{pdist} backward on zero rows: integer division by zero kills the process & Eg & C & P & C & F \\
B17 & \code{random.seed} inside the function is ignored on the first call & Dy & S & Q & V & F \\
B18 & named tuple \code{==} with tensor fields: \code{RecursionError} (Python 3.14 only) & Dy & F & Q & E & P \\
B19 & \code{bmm} under autotuning: missing symbol export in a C++ template (Windows) & In & F & Q & E & P \\
B20 & AOTInductor cannot package models returning \code{torch.return\_types.*} & Ex & F & Q & E & F \\
B21 & \code{lerp} of Boolean tensors with a 0-d weight fails to compile & De & F & P & E & F \\
B22 & \code{sum} / \code{prod} with \code{dtype=torch.bool}: C++ error / assertion & In & F & P & E & P \\
B23 & \code{addmm} with a 0-d bias under autotuning: \code{IndexError} in lowering & In & F & Q & E & F \\
B24 & \code{var\_mean} / \code{std\_mean} of an empty tensor: mean is 0 instead of NaN & In & S & P & V & F \\
B25 & \code{channel\_shuffle} loses the \code{channels\_last} layout & In & S & P & D & C \\
B26 & \code{str(KeyError(k))} loses the quotes that CPython adds & Dy & S & P & V & F \\
B27 & \code{abs(complex, out=float64)} rejected by the meta kernel & De & F & P & E & C \\
C1  & in-place-modification check lost for a saved input view: silent wrong gradients & AOT & S & P & G & F \\
C2  & \code{adaptive\_max\_pool3d} backward wrong for \code{channels\_last\_3d} (eager kernel) & Eg & S & P & G & C \\
C3  & a changing Python float argument is baked into the cached graph & Dy & S & Q & V & P \\
C7  & Numba: \code{remainder} of \code{MIN\_INT} by $-1$ kills the process & Ext & C & -- & C & P \\
C8  & \code{x * 1} / \code{x + 0} as a graph output return the input object itself & In & S & P & I & C \\
C9  & \code{torch.fx} codegen: \code{operator.pow(-2, x)} emitted as \code{-2**x} & Fx & S & P & V & F \\
C10 & floating-point floor division off by one (\code{1.0 // 0.1} gives 10) & In & S & P & V & C \\
C11 & closure mutations in methods of a locally defined class are dropped & Dy & S & P & S & F \\
C12 & JAX: \code{gcd} / \code{lcm} hang on the integer minimum & Ext & C & -- & C & F \\
C13 & compiled NumPy code: \code{fix}, \code{cbrt}, \code{clip}, \code{sign} return wrong values & Dy & S & P & V & P \\
C14 & vectorized \code{remainder}: NaN for an infinite divisor, 0 for large quotients & In & S & P & V & P \\
C16 & \code{round(SymInt, -1)} is the identity & Dy & S & P & V & P \\
C17 & NaN self-comparison of an \code{.item()} value folded to a constant & Dy & S & P & V & P \\
C20 & compiled NumPy code: dtype and shape rules differ (\code{cumsum}, \code{square}, \code{median}) & Dy & S & P & D & P \\
C22 & \code{clamp(int8, -1000, 1000)} returns $-24$: bounds wrapped into int8 & In & S & P & V & P \\
C24 & \code{torch.fx}: parameters named \code{nan} / \code{inf} shadow float constants & Fx & S & P & V & P \\
C27 & \code{GraphModule} pickle round trip turns in-place nodes into out-of-place ones & Fx & S & Q & S & P \\
C29 & \code{any(uint8)} returns \code{bool} instead of \code{uint8} & In & S & P & D & P \\
C32 & complex \code{pinv} / \code{polar} / \code{matrix\_sqrth} wrong under any dispatch mode & DM & S & Q & V & P \\
C34 & \code{avg\_pool} backward with \code{ceil\_mode}: full kernel size as divisor for the last window & In & S & P & G & P \\
C35 & \code{erfinv} loses four digits near $\pm 1$ & In & S & P & R & P \\
C36 & vectorized float64 \code{acosh} overflows above $1.3\times10^{154}$ & In & S & P & V & P \\
C37 & 134 operators compute on dtypes the eager kernel rejects (418 pairs) & In & S & P & E & P \\
C40 & vectorized float32 \code{erf} returns 0 below $10^{-7}$; \code{atanh} returns 0 below $\epsilon$ & In & S & P & R & P \\
C41 & \code{OrderedDict.move\_to\_end} on an input dictionary is dropped & Dy & S & Q & S & P \\
C42 & a generator that leaves the compiled region returns as an exhausted tuple iterator & Dy & S & Q & I & P \\
C44 & \code{IndexError} from a tensor operator bypasses the program's \code{except} clause & Dy & S & P & E & P \\
C46 & AOTInductor runner destroyed after a run: access violation (Windows, OpenMP workers still spinning) & Ex & C & Q & C & P \\
\bottomrule
\end{tabular}
\end{table}


\paragraph{Findings over time.} The first 27 issues (B1--B27) were found with default paths; 14 of them are
failures of programs that eager execution accepts. The issues found after the path factors and the
metadata-derived factors were added (C1--C46) are silent in 26 of 28 cases and include the four Dynamo defects
of RQ1, the identity defect C8, the 134-operator family C37, and the three accuracy defects. The maintainers'
own umbrella issue on lost validation, opened during the study, lists several of our reports as members.

%% Figure: (a) heatmap factor group x layer over the 68 PyTorch items, (b) stacked bars symptom per layer.
\newcommand{\hcell}[3]{% x, y, count
  \pgfmathsetmacro{\shade}{min(100, 12 + 7*#3)}
  \ifnum#3>0
    \fill[tolblue!\shade!white] (#1-0.5,#2-0.5) rectangle (#1+0.5,#2+0.5);
    \node[font=\scriptsize, text=\ifnum#3>7 white\else black\fi] at (#1,#2) {#3};
  \else
    \node[font=\scriptsize, black!35] at (#1,#2) {$\cdot$};
  \fi}
\begin{figure}[t]
\centering
\begin{minipage}[b]{0.58\linewidth}
\centering
\begin{tikzpicture}[x=0.66cm, y=0.42cm]
  %% columns: 1 Dy, 2 AOT, 3 De, 4 In, 5 Ex, 6 Eg, 7 Fx, 8 DM ; rows top->bottom 8..1
  \foreach \x/\lab in {1/Dy, 2/AOT, 3/De, 4/In, 5/Ex, 6/Eg, 7/Fx, 8/DM} \node[font=\scriptsize] at (\x,8.9) {\lab};
  \foreach \y/\lab in {8/{Q: graph-break position}, 7/{Q: backend / context}, 6/{Q: artifact route},
                       5/{Q: context sequence}, 4/{P: boundary values}, 3/{P: dtype factors},
                       2/{P: source factors}, 1/{P: Python corpus}}
    \node[font=\scriptsize, anchor=east] at (0.4,\y) {\lab};
  %% data (findings.csv, items with a PyTorch layer; cross-target items excluded)
  \hcell{1}{8}{3}\hcell{2}{8}{0}\hcell{3}{8}{0}\hcell{4}{8}{0}\hcell{5}{8}{0}\hcell{6}{8}{0}\hcell{7}{8}{0}\hcell{8}{8}{0}
  \hcell{1}{7}{3}\hcell{2}{7}{2}\hcell{3}{7}{0}\hcell{4}{7}{5}\hcell{5}{7}{0}\hcell{6}{7}{0}\hcell{7}{7}{0}\hcell{8}{7}{3}
  \hcell{1}{6}{0}\hcell{2}{6}{0}\hcell{3}{6}{0}\hcell{4}{6}{0}\hcell{5}{6}{3}\hcell{6}{6}{0}\hcell{7}{6}{1}\hcell{8}{6}{0}
  \hcell{1}{5}{4}\hcell{2}{5}{0}\hcell{3}{5}{0}\hcell{4}{5}{0}\hcell{5}{5}{0}\hcell{6}{5}{0}\hcell{7}{5}{0}\hcell{8}{5}{0}
  \hcell{1}{4}{1}\hcell{2}{4}{0}\hcell{3}{4}{3}\hcell{4}{4}{13}\hcell{5}{4}{0}\hcell{6}{4}{2}\hcell{7}{4}{0}\hcell{8}{4}{0}
  \hcell{1}{3}{1}\hcell{2}{3}{1}\hcell{3}{3}{2}\hcell{4}{3}{10}\hcell{5}{3}{0}\hcell{6}{3}{0}\hcell{7}{3}{0}\hcell{8}{3}{0}
  \hcell{1}{2}{2}\hcell{2}{2}{1}\hcell{3}{2}{0}\hcell{4}{2}{1}\hcell{5}{2}{0}\hcell{6}{2}{1}\hcell{7}{2}{0}\hcell{8}{2}{0}
  \hcell{1}{1}{4}\hcell{2}{1}{0}\hcell{3}{1}{0}\hcell{4}{1}{0}\hcell{5}{1}{0}\hcell{6}{1}{0}\hcell{7}{1}{2}\hcell{8}{1}{0}
  \draw[black!40] (0.5,0.5) rectangle (8.5,8.5);
  \foreach \x in {1.5,2.5,...,7.5} \draw[black!12] (\x,0.5) -- (\x,8.5);
  \foreach \y in {1.5,2.5,...,7.5} \draw[black!12] (0.5,\y) -- (8.5,\y);
  \draw[black!60, thick] (0.5,4.5) -- (8.5,4.5);
  %% marginals
  \foreach \x/\s in {1/18, 2/4, 3/5, 4/29, 5/3, 6/3, 7/3, 8/3} \node[font=\scriptsize\bfseries] at (\x,-0.1) {\s};
  \node[font=\scriptsize, anchor=east] at (0.4,-0.1) {items per layer};
\end{tikzpicture}
\subcaption{Items by factor group (Q = path side, P = program side) and layer.}
\end{minipage}\hfill
\begin{minipage}[b]{0.40\linewidth}
\centering
\begin{tikzpicture}
\begin{axis}[
    tccaxis, width=\linewidth, height=4.55cm,
    xbar stacked, bar width=5.5pt,
    ytick={1,2,3,4,5,6,7,8}, yticklabels={DM, Fx, Eg, Ex, In, De, AOT, Dy},
    xmin=0, xmax=32, xtick={0,10,20,30}, xlabel={items}, ymajorgrids=false, xmajorgrids,
    legend style={at={(0.98,0.03)}, anchor=south east}, legend columns=1,
]
  \addplot[fill=tolred!65, draw=tolred!80!black] coordinates {(3,1) (3,2) (1,3) (0,4) (24,5) (2,6) (2,7) (15,8)};
  \addplot[fill=tolblue!40, draw=tolblue!80!black] coordinates {(0,1) (0,2) (0,3) (2,4) (5,5) (3,6) (1,7) (3,8)};
  \addplot[fill=black!70, draw=black] coordinates {(0,1) (0,2) (2,3) (1,4) (0,5) (0,6) (1,7) (0,8)};
  \legend{silent (50), failure (14), crash (4)}
\end{axis}
\end{tikzpicture}
\subcaption{Symptoms per layer.}
\end{minipage}
\caption{RQ6: distribution of the 68 PyTorch items (issues and comments) by factor group, layer, and symptom.
Layer abbreviations as in Table~\ref{tab:bugs}; the five cross-compiler items are omitted.}
\label{fig:rq5}
\end{figure}


\paragraph{Cost and negative results.} The compile-based campaigns on default paths were the most expensive and the least productive:
about 2{,}800 programs and 45{,}000 differential tests on default paths produced one new defect, the
fabricated mean of Figure~\ref{fig:motivation}b. The path-side and metadata-derived analyzers produced the
remaining defects at a fraction of the cost. Table~\ref{tab:negative} lists the factor families that found
nothing new; they bound the claim of this paper and tell developers where consistency already holds.


\rqanswer{RQ6}{\tool led to 73 reports against PyTorch~2.14 and nightly builds and two other compilers; 14
are fixed or have a fix under review and 10 are confirmed. Of the 58 pipeline findings 42 are silent, and none
appears with a contiguous float32 tensor on the default path.}
